\chapter*{How to use this encyclopaedia}
\addcontentsline{toc}{chapter}{How to use this encyclopaedia}
% STATUS: written 2026-11-06 (describes the organisation the book already has:
% the eight question families of Section~\ref{sec:basic-theory:question-families},
% the chapter layout, and the two appendices. No new claim is made.)
\label{page:how-to-use}

The book is organised around eight questions, not around its chapters.
A reader who wants one of them answered can go straight to the
appendix that indexes it. A reader who wants to follow the subject from its
beginnings can read Part~I onward.

\section*{The eight question families}

Section~\ref{sec:basic-theory:question-families} fixes the eight questions the
book asks of every class of algebra. They are:

\begin{enumerate}
\item Which kind of derivation implies which other?
\item When are two definitions equivalent?
\item Is every derivation inner?
\item Is every derivation continuous, and does continuity come for free?
\item Does the algebra have to be commutative?
\item Can all derivations of the class be described?
\item Does a derivation extend or lift to a larger algebra?
\item Where does an expected implication fail?
\end{enumerate}

The chapters answer them in this order wherever the sources allow.
Appendix~\ref{app:questions} collects the answers in tables, family by family,
so that a reader comparing two classes of algebra can see both at once.

\section*{How a chapter is laid out}

Each chapter follows the same plan.

The introduction gives the history: why the topic arose, who first asked the
question, and what the chapter answers.
The sections that follow begin with the definitions as their sources state them,
then the theorems as their sources state them, each with the page or numbered
item.

Counterexamples come where a family of the last kind fails; they are what
answers family~8.
An application comes where the notion it uses is defined, so the reader meets
the use at the same point as the idea.
An open problem comes where the chapter states the question it belongs to,
and Appendix~\ref{app:questions} gathers all of them by family and by class.

A recalled theorem keeps the kind of its source: a proposition stays a
proposition, and a theorem stays a theorem.
The book never states a recalled result again in its own setting.
The specialisation is written after the recall, and it refers to the recall by
number.

\section*{The two appendix chapters}

Appendix~\ref{app:atlas} is the atlas of the literature.
It gives the shape of the bibliography: the works by decade, by language, and by
the kind of derivation-like map they treat.
It is counted again with each edition, so it records what this edition has read
rather than what the field contains.

Appendix~\ref{app:questions} is the collection of question tables.
For each of the eight families it gives, class by class, whether the answer is
known, false with a counterexample, or open, with the source that settles or
asks it.

\section*{Reading paths}

A reader new to the subject can read Part~I in order; it carries the history and
the basic theory, and it fixes the question families.
A reader who works on rings can take Part~II and Part~III, which treat the
associative classes, and consult Part~VI for the analytic ones.
A reader who comes from analysis or operator theory can start with Part~VI and
use Part~II for the algebraic background.
A reader who wants the current frontier can read
Appendix~\ref{app:questions}, which names every open problem with the chapter
that states it, then follow those references back.

\section*{The living edition}

The edition is updated each month.
New work is added where it changes a statement, a table entry or an open
problem.
A work that only confirms what the book already says is not cited for that
alone.
The version number appears on the title page, so that a citation to this book
can name the edition it used.
